% Options for packages loaded elsewhere
\PassOptionsToPackage{unicode}{hyperref}
\PassOptionsToPackage{hyphens}{url}
\PassOptionsToPackage{dvipsnames,svgnames,x11names}{xcolor}
%
\documentclass[
  a4paper,
]{article}

\usepackage{amsmath,amssymb}
\usepackage{iftex}
\ifPDFTeX
  \usepackage[T1]{fontenc}
  \usepackage[utf8]{inputenc}
  \usepackage{textcomp} % provide euro and other symbols
\else % if luatex or xetex
  \usepackage{unicode-math}
  \defaultfontfeatures{Scale=MatchLowercase}
  \defaultfontfeatures[\rmfamily]{Ligatures=TeX,Scale=1}
\fi
\usepackage{lmodern}
\ifPDFTeX\else  
    % xetex/luatex font selection
\fi
% Use upquote if available, for straight quotes in verbatim environments
\IfFileExists{upquote.sty}{\usepackage{upquote}}{}
\IfFileExists{microtype.sty}{% use microtype if available
  \usepackage[]{microtype}
  \UseMicrotypeSet[protrusion]{basicmath} % disable protrusion for tt fonts
}{}
\makeatletter
\@ifundefined{KOMAClassName}{% if non-KOMA class
  \IfFileExists{parskip.sty}{%
    \usepackage{parskip}
  }{% else
    \setlength{\parindent}{0pt}
    \setlength{\parskip}{6pt plus 2pt minus 1pt}}
}{% if KOMA class
  \KOMAoptions{parskip=half}}
\makeatother
\usepackage{xcolor}
\usepackage[margin=2cm,top=2.5cm,bottom=2.5cm]{geometry}
\setlength{\emergencystretch}{3em} % prevent overfull lines
\setcounter{secnumdepth}{5}
% Make \paragraph and \subparagraph free-standing
\makeatletter
\ifx\paragraph\undefined\else
  \let\oldparagraph\paragraph
  \renewcommand{\paragraph}{
    \@ifstar
      \xxxParagraphStar
      \xxxParagraphNoStar
  }
  \newcommand{\xxxParagraphStar}[1]{\oldparagraph*{#1}\mbox{}}
  \newcommand{\xxxParagraphNoStar}[1]{\oldparagraph{#1}\mbox{}}
\fi
\ifx\subparagraph\undefined\else
  \let\oldsubparagraph\subparagraph
  \renewcommand{\subparagraph}{
    \@ifstar
      \xxxSubParagraphStar
      \xxxSubParagraphNoStar
  }
  \newcommand{\xxxSubParagraphStar}[1]{\oldsubparagraph*{#1}\mbox{}}
  \newcommand{\xxxSubParagraphNoStar}[1]{\oldsubparagraph{#1}\mbox{}}
\fi
\makeatother

\usepackage{color}
\usepackage{fancyvrb}
\newcommand{\VerbBar}{|}
\newcommand{\VERB}{\Verb[commandchars=\\\{\}]}
\DefineVerbatimEnvironment{Highlighting}{Verbatim}{commandchars=\\\{\}}
% Add ',fontsize=\small' for more characters per line
\usepackage{framed}
\definecolor{shadecolor}{RGB}{241,243,245}
\newenvironment{Shaded}{\begin{snugshade}}{\end{snugshade}}
\newcommand{\AlertTok}[1]{\textcolor[rgb]{0.68,0.00,0.00}{#1}}
\newcommand{\AnnotationTok}[1]{\textcolor[rgb]{0.37,0.37,0.37}{#1}}
\newcommand{\AttributeTok}[1]{\textcolor[rgb]{0.40,0.45,0.13}{#1}}
\newcommand{\BaseNTok}[1]{\textcolor[rgb]{0.68,0.00,0.00}{#1}}
\newcommand{\BuiltInTok}[1]{\textcolor[rgb]{0.00,0.23,0.31}{#1}}
\newcommand{\CharTok}[1]{\textcolor[rgb]{0.13,0.47,0.30}{#1}}
\newcommand{\CommentTok}[1]{\textcolor[rgb]{0.37,0.37,0.37}{#1}}
\newcommand{\CommentVarTok}[1]{\textcolor[rgb]{0.37,0.37,0.37}{\textit{#1}}}
\newcommand{\ConstantTok}[1]{\textcolor[rgb]{0.56,0.35,0.01}{#1}}
\newcommand{\ControlFlowTok}[1]{\textcolor[rgb]{0.00,0.23,0.31}{\textbf{#1}}}
\newcommand{\DataTypeTok}[1]{\textcolor[rgb]{0.68,0.00,0.00}{#1}}
\newcommand{\DecValTok}[1]{\textcolor[rgb]{0.68,0.00,0.00}{#1}}
\newcommand{\DocumentationTok}[1]{\textcolor[rgb]{0.37,0.37,0.37}{\textit{#1}}}
\newcommand{\ErrorTok}[1]{\textcolor[rgb]{0.68,0.00,0.00}{#1}}
\newcommand{\ExtensionTok}[1]{\textcolor[rgb]{0.00,0.23,0.31}{#1}}
\newcommand{\FloatTok}[1]{\textcolor[rgb]{0.68,0.00,0.00}{#1}}
\newcommand{\FunctionTok}[1]{\textcolor[rgb]{0.28,0.35,0.67}{#1}}
\newcommand{\ImportTok}[1]{\textcolor[rgb]{0.00,0.46,0.62}{#1}}
\newcommand{\InformationTok}[1]{\textcolor[rgb]{0.37,0.37,0.37}{#1}}
\newcommand{\KeywordTok}[1]{\textcolor[rgb]{0.00,0.23,0.31}{\textbf{#1}}}
\newcommand{\NormalTok}[1]{\textcolor[rgb]{0.00,0.23,0.31}{#1}}
\newcommand{\OperatorTok}[1]{\textcolor[rgb]{0.37,0.37,0.37}{#1}}
\newcommand{\OtherTok}[1]{\textcolor[rgb]{0.00,0.23,0.31}{#1}}
\newcommand{\PreprocessorTok}[1]{\textcolor[rgb]{0.68,0.00,0.00}{#1}}
\newcommand{\RegionMarkerTok}[1]{\textcolor[rgb]{0.00,0.23,0.31}{#1}}
\newcommand{\SpecialCharTok}[1]{\textcolor[rgb]{0.37,0.37,0.37}{#1}}
\newcommand{\SpecialStringTok}[1]{\textcolor[rgb]{0.13,0.47,0.30}{#1}}
\newcommand{\StringTok}[1]{\textcolor[rgb]{0.13,0.47,0.30}{#1}}
\newcommand{\VariableTok}[1]{\textcolor[rgb]{0.07,0.07,0.07}{#1}}
\newcommand{\VerbatimStringTok}[1]{\textcolor[rgb]{0.13,0.47,0.30}{#1}}
\newcommand{\WarningTok}[1]{\textcolor[rgb]{0.37,0.37,0.37}{\textit{#1}}}

\providecommand{\tightlist}{%
  \setlength{\itemsep}{0pt}\setlength{\parskip}{0pt}}\usepackage{longtable,booktabs,array}
\usepackage{calc} % for calculating minipage widths
% Correct order of tables after \paragraph or \subparagraph
\usepackage{etoolbox}
\makeatletter
\patchcmd\longtable{\par}{\if@noskipsec\mbox{}\fi\par}{}{}
\makeatother
% Allow footnotes in longtable head/foot
\IfFileExists{footnotehyper.sty}{\usepackage{footnotehyper}}{\usepackage{footnote}}
\makesavenoteenv{longtable}
\usepackage{graphicx}
\makeatletter
\newsavebox\pandoc@box
\newcommand*\pandocbounded[1]{% scales image to fit in text height/width
  \sbox\pandoc@box{#1}%
  \Gscale@div\@tempa{\textheight}{\dimexpr\ht\pandoc@box+\dp\pandoc@box\relax}%
  \Gscale@div\@tempb{\linewidth}{\wd\pandoc@box}%
  \ifdim\@tempb\p@<\@tempa\p@\let\@tempa\@tempb\fi% select the smaller of both
  \ifdim\@tempa\p@<\p@\scalebox{\@tempa}{\usebox\pandoc@box}%
  \else\usebox{\pandoc@box}%
  \fi%
}
% Set default figure placement to htbp
\def\fps@figure{htbp}
\makeatother

\makeatletter
\@ifpackageloaded{caption}{}{\usepackage{caption}}
\AtBeginDocument{%
\ifdefined\contentsname
  \renewcommand*\contentsname{Table des matières}
\else
  \newcommand\contentsname{Table des matières}
\fi
\ifdefined\listfigurename
  \renewcommand*\listfigurename{Liste des Figures}
\else
  \newcommand\listfigurename{Liste des Figures}
\fi
\ifdefined\listtablename
  \renewcommand*\listtablename{Liste des Tables}
\else
  \newcommand\listtablename{Liste des Tables}
\fi
\ifdefined\figurename
  \renewcommand*\figurename{Figure}
\else
  \newcommand\figurename{Figure}
\fi
\ifdefined\tablename
  \renewcommand*\tablename{Table}
\else
  \newcommand\tablename{Table}
\fi
}
\@ifpackageloaded{float}{}{\usepackage{float}}
\floatstyle{ruled}
\@ifundefined{c@chapter}{\newfloat{codelisting}{h}{lop}}{\newfloat{codelisting}{h}{lop}[chapter]}
\floatname{codelisting}{Listing}
\newcommand*\listoflistings{\listof{codelisting}{Liste des Listings}}
\makeatother
\makeatletter
\makeatother
\makeatletter
\@ifpackageloaded{caption}{}{\usepackage{caption}}
\@ifpackageloaded{subcaption}{}{\usepackage{subcaption}}
\makeatother

\ifLuaTeX
\usepackage[bidi=basic]{babel}
\else
\usepackage[bidi=default]{babel}
\fi
\babelprovide[main,import]{french}
% get rid of language-specific shorthands (see #6817):
\let\LanguageShortHands\languageshorthands
\def\languageshorthands#1{}
\usepackage{bookmark}

\IfFileExists{xurl.sty}{\usepackage{xurl}}{} % add URL line breaks if available
\urlstyle{same} % disable monospaced font for URLs
\hypersetup{
  pdftitle={PostgreSQL sur Windows --- Guide d'Installation},
  pdfauthor={Guide pour débutants},
  pdflang={fr},
  colorlinks=true,
  linkcolor={blue},
  filecolor={Maroon},
  citecolor={Blue},
  urlcolor={Blue},
  pdfcreator={LaTeX via pandoc}}


\title{PostgreSQL sur Windows --- Guide d'Installation}
\author{Guide pour débutants}
\date{2025-10-27}

\begin{document}
\maketitle

\renewcommand*\contentsname{Table des matières}
{
\hypersetup{linkcolor=}
\setcounter{tocdepth}{2}
\tableofcontents
}

Ce guide vous accompagne étape par étape : installation de PostgreSQL,
vérification de votre configuration, utilisation de la console
\texttt{psql}, chargement de la base de données exemple Pagila, et
utilisation de clients GUI conviviaux pour débutants. Des conseils
clairs et des solutions rapides sont inclus tout au long du guide.

\begin{center}\rule{0.5\linewidth}{0.5pt}\end{center}

\section{Avant de Commencer}\label{avant-de-commencer}

\textbf{Vous aurez besoin de} - Un PC Windows avec droits administrateur
et quelques Go d'espace disque libre. - Un navigateur web moderne pour
télécharger les fichiers. - Environ \textbf{30--60 minutes} dans une
session calme et sans interruption.

\textbf{Termes clés (rapide)} - \textbf{Server} --- le moteur de base de
données PostgreSQL fonctionnant sur votre PC. - \textbf{Client} --- les
outils que vous utilisez pour communiquer avec le server (ex. :
\texttt{psql}, pgAdmin, DBeaver). - \textbf{Database} --- un conteneur
de tables, vues, etc. (un seul server peut héberger plusieurs
databases). - \textbf{Role/User} --- un compte qui se connecte et
possède des permissions.

\begin{center}\rule{0.5\linewidth}{0.5pt}\end{center}

\section{Installer PostgreSQL
(Windows)}\label{installer-postgresql-windows}

\subsection{Télécharger PostgreSQL}\label{tuxe9luxe9charger-postgresql}

\textbf{Lien officiel :}
\url{https://www.postgresql.org/download/windows/}

\begin{enumerate}
\def\labelenumi{\arabic{enumi}.}
\tightlist
\item
  Ouvrez la page \textbf{officielle PostgreSQL pour Windows} et
  choisissez \textbf{Interactive installer by EDB}.
\item
  Téléchargez la dernière version stable pour \textbf{Windows x86‑64}.

  \begin{itemize}
  \tightlist
  \item
    \textbf{Lien direct :}
    \url{https://www.enterprisedb.com/downloads/postgres-postgresql-downloads}
  \end{itemize}
\item
  Exécutez l'installateur (\texttt{.exe}) et conservez les composants
  par défaut : \textbf{Server}, \textbf{pgAdmin}, \textbf{Command Line
  Tools}.
\item
  Lorsqu'on vous le demande, \textbf{définissez un mot de passe pour le
  superuser \texttt{postgres}} et conservez-le en lieu sûr.
\item
  Conservez le \textbf{Port} à \texttt{5432} et acceptez la
  \textbf{Locale} par défaut.
\item
  Terminez l'assistant d'installation.
\end{enumerate}

\begin{quote}
\textbf{Composants installés} - \textbf{PostgreSQL Server} --- le moteur
qui stocke vos données\\
- \textbf{\texttt{psql}} --- le shell SQL en ligne de commande\\
- \textbf{pgAdmin} --- l'outil de gestion graphique officiel
\end{quote}

\subsection{Vérifier l'installation (au
choix)}\label{vuxe9rifier-linstallation-au-choix}

\textbf{Option A --- Menu Démarrer} - Menu Démarrer → \textbf{SQL Shell
(psql)} → appuyez sur \textbf{Entrée} pour chaque invite → entrez votre
mot de passe. Si vous voyez \texttt{postgres=\#}, c'est bon.

\textbf{Option B --- Invite de commandes (CMD)}

\begin{Shaded}
\begin{Highlighting}[]
\StringTok{"C:\textbackslash{}Program Files\textbackslash{}PostgreSQL\textbackslash{}17\textbackslash{}bin\textbackslash{}psql.exe"}\NormalTok{ {-}{-}version}
\end{Highlighting}
\end{Shaded}

\emph{(remplacez \texttt{17} par votre version installée)}

\textbf{Sortie attendue :}

\begin{verbatim}
psql (PostgreSQL) 17.0
\end{verbatim}

\textbf{Si Windows indique ``psql n'est pas reconnu''}

\begin{itemize}
\tightlist
\item
  Utilisez \textbf{Démarrer → SQL Shell (psql)} (fonctionne sans PATH),
  \textbf{ou}
\item
  Ajoutez le dossier bin de PostgreSQL à votre \textbf{PATH} :

  \begin{itemize}
  \tightlist
  \item
    Système → À propos → Paramètres système avancés
  \item
    Variables d'environnement → Path → Modifier → Nouveau
  \item
    Ajoutez :
    \texttt{C:\textbackslash{}Program\ Files\textbackslash{}PostgreSQL\textbackslash{}17\textbackslash{}bin}
  \item
    Puis rouvrez votre terminal
  \end{itemize}
\end{itemize}

\begin{center}\rule{0.5\linewidth}{0.5pt}\end{center}

\section{Démarrer/Arrêter le Service de Base de Données (si
nécessaire)}\label{duxe9marrerarruxeater-le-service-de-base-de-donnuxe9es-si-nuxe9cessaire}

\begin{itemize}
\tightlist
\item
  Appuyez sur \textbf{Win + R}, tapez \texttt{services.msc}, appuyez sur
  \textbf{Entrée}.
\item
  Trouvez \textbf{PostgreSQL }. Clic droit → \textbf{Démarrer} (ou
  \textbf{Redémarrer/Arrêter}). Il démarre normalement automatiquement
  avec Windows.
\end{itemize}

Si vous voyez des timeouts de connexion ou ``could not connect to
server'', vérifiez d'abord ce service.

\begin{center}\rule{0.5\linewidth}{0.5pt}\end{center}

\section{\texorpdfstring{Première Connexion avec
\texttt{psql}}{Première Connexion avec psql}}\label{premiuxe8re-connexion-avec-psql}

\begin{enumerate}
\def\labelenumi{\arabic{enumi}.}
\tightlist
\item
  Ouvrez \textbf{SQL Shell (psql)} depuis le menu Démarrer.
\item
  Appuyez sur \textbf{Entrée} à chaque invite pour accepter les valeurs
  par défaut (Server \texttt{localhost}, Database \texttt{postgres},
  Port \texttt{5432}, User \texttt{postgres}).
\item
  Entrez votre \textbf{mot de passe \texttt{postgres}}. Vous devriez
  voir une invite comme \texttt{postgres=\#}.
\end{enumerate}

\subsection{\texorpdfstring{Qu'est-ce que \texttt{psql} ? (introduction
douce)}{Qu'est-ce que psql ? (introduction douce)}}\label{quest-ce-que-psql-introduction-douce}

\begin{itemize}
\tightlist
\item
  \textbf{\texttt{psql}} est une console textuelle pour communiquer avec
  PostgreSQL. Vous tapez une commande, appuyez sur \textbf{Entrée}, et
  le server répond.
\item
  \textbf{Pourquoi l'apprendre ?} C'est rapide, correspond à la
  documentation officielle, et vous aide à résoudre des problèmes même
  quand un GUI ne peut pas se connecter.
\item
  \textbf{Ce que vous verrez :} une invite comme \texttt{postgres=\#}
  (ou le nom de votre database). Terminez les instructions SQL par un
  \textbf{point-virgule} (\texttt{;}). Si vous l'oubliez, \texttt{psql}
  attendra plus d'entrée.
\item
  \textbf{Deux types de commandes :}

  \begin{itemize}
  \tightlist
  \item
    \textbf{SQL} (se termine par \texttt{;}) --- ex. :
    \texttt{SELECT\ version();}
  \item
    \textbf{Meta‑commands} (commencent par un backslash) --- ex. :
    \texttt{\textbackslash{}dt} (lister les tables),
    \texttt{\textbackslash{}d\ film} (décrire une table). Ce sont des
    fonctionnalités de \texttt{psql} lui-même, pas du SQL.
  \end{itemize}
\item
  \textbf{Mini démonstration :}
\end{itemize}

\begin{Shaded}
\begin{Highlighting}[]
\KeywordTok{SELECT} \DecValTok{1}\NormalTok{;}
\KeywordTok{CREATE} \KeywordTok{DATABASE}\NormalTok{ testdb;}
\NormalTok{\textbackslash{}c testdb}
\KeywordTok{CREATE} \KeywordTok{TABLE}\NormalTok{ t(x }\DataTypeTok{int}\NormalTok{);}
\KeywordTok{INSERT} \KeywordTok{INTO}\NormalTok{ t }\KeywordTok{VALUES}\NormalTok{ (}\DecValTok{1}\NormalTok{),(}\DecValTok{2}\NormalTok{);}
\KeywordTok{SELECT} \OperatorTok{*} \KeywordTok{FROM}\NormalTok{ t;}
\end{Highlighting}
\end{Shaded}

\begin{itemize}
\tightlist
\item
  \textbf{Erreurs courantes des débutants :} point-virgule manquant ;
  connexion à la mauvaise database ; tout faire en tant que superuser
  \texttt{postgres} (créez un user normal pour le travail quotidien).
\end{itemize}

\begin{quote}
\textbf{Analogie :} PostgreSQL est la \textbf{bibliothèque} (server).
\texttt{psql} est le \textbf{bureau du bibliothécaire} : vous demandez
des choses (SQL) et obtenez des réponses. Les GUIs ajoutent des cartes
et des boutons, mais le bureau fonctionne toujours.
\end{quote}

\begin{center}\rule{0.5\linewidth}{0.5pt}\end{center}

\section{\texorpdfstring{Charger la Base de Données Exemple
\textbf{Pagila}}{Charger la Base de Données Exemple Pagila}}\label{charger-la-base-de-donnuxe9es-exemple-pagila}

Pagila est un exemple classique (magasin de location de DVD) avec des
tables réalistes (films, acteurs, clients, locations). C'est parfait
pour pratiquer les jointures, les agrégats et les clés étrangères.

\subsection{Télécharger les
fichiers}\label{tuxe9luxe9charger-les-fichiers}

\textbf{Lien officiel Pagila :}
\url{https://github.com/devrimgunduz/pagila}

Créez le dossier \texttt{C:\textbackslash{}Pagila\_Project} et
téléchargez les deux fichiers suivants :

\begin{itemize}
\tightlist
\item
  \textbf{Schema :}
  \href{https://raw.githubusercontent.com/devrimgunduz/pagila/master/pagila-schema.sql}{pagila-schema.sql}
\item
  \textbf{Data :}
  \href{https://raw.githubusercontent.com/devrimgunduz/pagila/master/pagila-insert-data.sql}{pagila-data.sql}
\end{itemize}

\emph{(Clic droit sur les liens → Enregistrer sous\ldots{} → Enregistrer
dans \texttt{C:\textbackslash{}Pagila\_Project})}

\subsection{Créer la database}\label{cruxe9er-la-database}

Dans \texttt{psql} :

\begin{Shaded}
\begin{Highlighting}[]
\KeywordTok{CREATE} \KeywordTok{DATABASE}\NormalTok{ pagila;}
\NormalTok{\textbackslash{}c pagila}
\end{Highlighting}
\end{Shaded}

\subsection{Charger le schema et les données (choisissez une
méthode)}\label{charger-le-schema-et-les-donnuxe9es-choisissez-une-muxe9thode}

\textbf{Méthode A --- depuis l'intérieur de \texttt{psql} (plus facile)}

\begin{quote}
\textbf{Note :} Utilisez des slashes normaux (/) dans les chemins
\texttt{psql}, même sous Windows.
\end{quote}

\begin{Shaded}
\begin{Highlighting}[]
\NormalTok{\textbackslash{}i }\StringTok{\textquotesingle{}C:/Pagila\_Project/pagila{-}schema.sql\textquotesingle{}}
\NormalTok{\textbackslash{}i }\StringTok{\textquotesingle{}C:/Pagila\_Project/pagila{-}data.sql\textquotesingle{}}
\end{Highlighting}
\end{Shaded}

\textbf{Méthode B --- depuis l'Invite de commandes (CMD) ou PowerShell}

\textbf{Pour CMD :}

\begin{Shaded}
\begin{Highlighting}[]
\BuiltInTok{cd}\NormalTok{ C:\textbackslash{}Pagila\_Project}
\BuiltInTok{set} \VariableTok{PSQL}\NormalTok{=}\StringTok{"C:\textbackslash{}Program Files\textbackslash{}PostgreSQL\textbackslash{}17\textbackslash{}bin\textbackslash{}psql.exe"}
\PreprocessorTok{\%}\VariableTok{PSQL}\PreprocessorTok{\%}\NormalTok{ {-}U postgres }\AttributeTok{{-}d}\NormalTok{ pagila }\AttributeTok{{-}f}\NormalTok{ pagila}\AttributeTok{{-}schema}\NormalTok{.sql}
\PreprocessorTok{\%}\VariableTok{PSQL}\PreprocessorTok{\%}\NormalTok{ {-}U postgres }\AttributeTok{{-}d}\NormalTok{ pagila }\AttributeTok{{-}f}\NormalTok{ pagila}\AttributeTok{{-}data}\NormalTok{.sql}
\end{Highlighting}
\end{Shaded}

\textbf{Pour PowerShell :}

\begin{Shaded}
\begin{Highlighting}[]
\FunctionTok{cd}\NormalTok{ C}\OperatorTok{:}\NormalTok{\textbackslash{}Pagila\_Project}
\VariableTok{$PSQL} \OperatorTok{=} \StringTok{"C:\textbackslash{}Program Files\textbackslash{}PostgreSQL\textbackslash{}17\textbackslash{}bin\textbackslash{}psql.exe"}
\OperatorTok{\&} \VariableTok{$PSQL} \OperatorTok{{-}}\NormalTok{U postgres }\OperatorTok{{-}}\NormalTok{d pagila }\OperatorTok{{-}}\NormalTok{f pagila{-}schema}\OperatorTok{.}\FunctionTok{sql}
\OperatorTok{\&} \VariableTok{$PSQL} \OperatorTok{{-}}\NormalTok{U postgres }\OperatorTok{{-}}\NormalTok{d pagila }\OperatorTok{{-}}\NormalTok{f pagila{-}data}\OperatorTok{.}\FunctionTok{sql}
\end{Highlighting}
\end{Shaded}

\emph{(Remplacez \texttt{17} par votre version installée.)}

\subsection{Confirmer que ça a
fonctionné}\label{confirmer-que-uxe7a-a-fonctionnuxe9}

\begin{Shaded}
\begin{Highlighting}[]
\NormalTok{\textbackslash{}c pagila}
\NormalTok{\textbackslash{}dt}
\KeywordTok{SELECT} \FunctionTok{COUNT}\NormalTok{(}\OperatorTok{*}\NormalTok{) }\KeywordTok{FROM}\NormalTok{ film;       }\CommentTok{{-}{-} \textasciitilde{}1,000}
\KeywordTok{SELECT} \FunctionTok{COUNT}\NormalTok{(}\OperatorTok{*}\NormalTok{) }\KeywordTok{FROM}\NormalTok{ customer;   }\CommentTok{{-}{-} \textasciitilde{}599}
\end{Highlighting}
\end{Shaded}

\begin{center}\rule{0.5\linewidth}{0.5pt}\end{center}

\section{Vous Préférez un GUI ? Clients
Populaires}\label{vous-pruxe9fuxe9rez-un-gui-clients-populaires}

Choisissez \textbf{un} pour commencer ; vous pourrez toujours essayer
les autres plus tard. (Les débutants choisissent souvent
\textbf{pgAdmin} ou \textbf{DBeaver}.)

\subsection{A) pgAdmin (officiel)}\label{a-pgadmin-officiel}

\textbf{Télécharger :} \url{https://www.pgadmin.org/download/}

\begin{itemize}
\tightlist
\item
  \textbf{Idéal pour :} Administration + apprentissage des concepts
  PostgreSQL de base
\item
  \textbf{Connexion :} Ouvrir pgAdmin → Register → Server\ldots{} → Host
  \texttt{localhost}, Port \texttt{5432}, Database \texttt{postgres},
  User \texttt{postgres} (ou \texttt{student}), mot de passe →
  \textbf{Save}
\item
  \textbf{Utilisation :} Créer des databases/roles, naviguer, exécuter
  des requêtes dans \textbf{Query Tool}
\end{itemize}

\subsection{B) DBeaver Community (gratuit,
multi-DB)}\label{b-dbeaver-community-gratuit-multi-db}

\textbf{Télécharger :} \url{https://dbeaver.io/download/}

\begin{itemize}
\tightlist
\item
  \textbf{Idéal pour :} Requêtes quotidiennes et exploration de données
  ; fonctionne avec de nombreuses databases
\item
  \textbf{Connexion :} New Database Connection → PostgreSQL → Host
  \texttt{localhost}, Port \texttt{5432}, Database \texttt{pagila} (ou
  \texttt{postgres}), User + Password → \textbf{Test} → \textbf{Finish}
\item
  \textbf{Fonctionnalités sympas :} Diagrammes ER, éditeurs de données,
  import/export CSV/Excel
\end{itemize}

\subsection{C) Visual Studio Code + extension PostgreSQL
(gratuit)}\label{c-visual-studio-code-extension-postgresql-gratuit}

\textbf{Télécharger :}
\url{https://marketplace.visualstudio.com/items?itemName=ms-ossdata.vscode-postgresql}

\begin{itemize}
\tightlist
\item
  \textbf{Idéal pour :} Développeurs qui vivent dans VS Code
\item
  \textbf{Connexion :} Extensions → rechercher \textbf{PostgreSQL}
  (Microsoft) → Install → Command Palette → \textbf{PostgreSQL: Add
  Connection} → remplir host/port/database/user/password
\item
  \textbf{Fonctionnalités sympas :} Exécuter des requêtes à côté de
  votre code ; IntelliSense ; grille de résultats
\end{itemize}

\subsection{Quel GUI choisir ? (choix
rapides)}\label{quel-gui-choisir-choix-rapides}

\begin{itemize}
\tightlist
\item
  \textbf{Totalement nouveau aux databases ?} \textbf{pgAdmin}
  (officiel, choix sûr par défaut)
\item
  \textbf{Explorateur moderne et convivial ?} \textbf{DBeaver Community}
\item
  \textbf{Je vis dans VS Code} → Extension \textbf{PostgreSQL for VS
  Code}
\end{itemize}

\subsection{Comment se connecter (mêmes paramètres dans la plupart des
GUIs)}\label{comment-se-connecter-muxeames-paramuxe8tres-dans-la-plupart-des-guis}

\begin{itemize}
\tightlist
\item
  \textbf{Host:} \texttt{localhost}
\item
  \textbf{Port:} \texttt{5432} (sauf si vous l'avez changé)
\item
  \textbf{Database:} \texttt{pagila} (après création) ou
  \texttt{postgres} pour commencer
\item
  \textbf{User:} \texttt{postgres} (admin) ou votre propre user (ex. :
  \texttt{student})
\item
  \textbf{Password:} celui que vous avez défini lors de l'installation
  ou pour ce role
\end{itemize}

\subsection{Conseils pro}\label{conseils-pro}

\begin{itemize}
\tightlist
\item
  Enregistrez la connexion pour ne pas avoir à retaper les identifiants.
\item
  Créez une seconde connexion pour votre user normal (pas
  \texttt{postgres}) pour pratiquer le principe du moindre privilège.
\item
  Apprenez la fonction \textbf{export vers CSV/Excel} du client pour les
  devoirs et rapports de laboratoire.
\end{itemize}

\subsection{Pièges courants des GUIs}\label{piuxe8ges-courants-des-guis}

\begin{itemize}
\tightlist
\item
  \textbf{Service non démarré :} Démarrez \textbf{Services} →
  \emph{PostgreSQL } → \textbf{Démarrer/Redémarrer}.
\item
  \textbf{Mauvais host/port :} Host par défaut \textbf{localhost}, port
  \textbf{5432}. Si \emph{localhost} échoue, essayez \textbf{127.0.0.1}.
\item
  \textbf{Identifiants/Database incorrects :} Assurez-vous que le
  \textbf{user} et la \textbf{database} existent ; testez avec
  \texttt{psql} d'abord.
\item
  \textbf{Mot de passe enregistré/ancien :} Mettez à jour le mot de
  passe enregistré dans votre GUI.
\item
  \textbf{Invites SSL (local) :} Pour le travail local, SSL est
  généralement inutile --- définissez le mode sur \textbf{Prefer} (ou
  \textbf{Disable} si autorisé) sauf si votre configuration nécessite
  SSL.
\end{itemize}

\subsection{Dépanner une connexion GUI --- flux
rapide}\label{duxe9panner-une-connexion-gui-flux-rapide}

\begin{verbatim}
1) PostgreSQL fonctionne-t-il ? Services → PostgreSQL <version> devrait être En cours d'exécution.

2) Le server répond-il ?

   **CMD:**
\end{verbatim}

``\%PGBIN\%\pg\_isready.exe'' -h localhost -p 5432 ```

\textbf{PowerShell:}
\texttt{\&\ "\$PGBIN\textbackslash{}pg\_isready.exe"\ -h\ localhost\ -p\ 5432}

\begin{enumerate}
\def\labelenumi{\arabic{enumi})}
\setcounter{enumi}{2}
\item
  Pouvez-vous vous connecter avec psql ?

  \textbf{CMD:}

\begin{verbatim}
"%PGBIN%\psql.exe" -U postgres -d postgres -c "SELECT 1;"
\end{verbatim}

  \textbf{PowerShell:}

\begin{verbatim}
& "$PGBIN\psql.exe" -U postgres -d postgres -c "SELECT 1;"
\end{verbatim}
\item
  Le port écoute-t-il ? CMD: netstat -ano \textbar{} findstr 5432
  PowerShell: Get-NetTCPConnection -LocalPort 5432
\item
  Essayez des variations de host : 127.0.0.1 au lieu de localhost (IPv4
  vs IPv6).
\item
  Réinitialisez/créez un user (si nécessaire) : ALTER ROLE postgres
  PASSWORD `NewStrongPassword!';
\item
  Pare-feu/outils de sécurité : autorisez postgres.exe sur le port 5432
  (réseau privé).
\end{enumerate}

\begin{verbatim}

> **Connexions distantes (plus tard) :** Pour qu'un autre ordinateur se connecte à votre server, vous devrez ajuster `postgresql.conf` (`listen_addresses`) et `pg_hba.conf`. Pour ce guide débutant (même PC), aucun changement n'est requis.

### Commandes Windows utiles (copier/coller)

**Pour CMD (Invite de commandes) :**

Définissez d'abord le chemin :
```cmd
set PGBIN=C:\Program Files\PostgreSQL\17\bin
\end{verbatim}

Puis utilisez :

\begin{Shaded}
\begin{Highlighting}[]
\StringTok{"}\PreprocessorTok{\%}\VariableTok{PGBIN}\PreprocessorTok{\%}\StringTok{\textbackslash{}psql.exe"}\NormalTok{ {-}{-}version}

\StringTok{"}\PreprocessorTok{\%}\VariableTok{PGBIN}\PreprocessorTok{\%}\StringTok{\textbackslash{}pg\_isready.exe"}\NormalTok{ {-}h localhost }\AttributeTok{{-}p}\NormalTok{ 5432}

\StringTok{"}\PreprocessorTok{\%}\VariableTok{PGBIN}\PreprocessorTok{\%}\StringTok{\textbackslash{}psql.exe"}\NormalTok{ {-}U postgres }\AttributeTok{{-}d}\NormalTok{ postgres }
\NormalTok{{-}c }\StringTok{"SELECT current\_database();"}
\end{Highlighting}
\end{Shaded}

\textbf{Pour PowerShell :}

Définissez d'abord le chemin :

\begin{Shaded}
\begin{Highlighting}[]
\VariableTok{$PGBIN} \OperatorTok{=} \StringTok{"C:\textbackslash{}Program Files\textbackslash{}PostgreSQL\textbackslash{}17\textbackslash{}bin"}
\end{Highlighting}
\end{Shaded}

Puis utilisez :

\begin{Shaded}
\begin{Highlighting}[]
\OperatorTok{\&} \StringTok{"}\VariableTok{$PGBIN}\StringTok{\textbackslash{}psql.exe"} \OperatorTok{{-}{-}}\NormalTok{version}

\OperatorTok{\&} \StringTok{"}\VariableTok{$PGBIN}\StringTok{\textbackslash{}pg\_isready.exe"} \OperatorTok{{-}}\NormalTok{h localhost }\OperatorTok{{-}}\NormalTok{p }\DecValTok{5432}

\OperatorTok{\&} \StringTok{"}\VariableTok{$PGBIN}\StringTok{\textbackslash{}psql.exe"} \OperatorTok{{-}}\NormalTok{U postgres }\OperatorTok{{-}}\NormalTok{d postgres }
\OperatorTok{{-}}\NormalTok{c }\StringTok{"SELECT current\_database();"}
\end{Highlighting}
\end{Shaded}

\begin{center}\rule{0.5\linewidth}{0.5pt}\end{center}

\section{\texorpdfstring{Essentiels \texttt{psql}
(Aide-Mémoire)}{Essentiels psql (Aide-Mémoire)}}\label{essentiels-psql-aide-muxe9moire}

\begin{Shaded}
\begin{Highlighting}[]
\NormalTok{\textbackslash{}q                 }\CommentTok{{-}{-} quitter psql}
\NormalTok{\textbackslash{}l                 }\CommentTok{{-}{-} lister les databases}
\NormalTok{\textbackslash{}c dbname          }\CommentTok{{-}{-} se connecter à une database (ex. : \textbackslash{}c pagila)}
\NormalTok{\textbackslash{}dt                }\CommentTok{{-}{-} lister les tables dans le schema actuel}
\NormalTok{\textbackslash{}d table\_name      }\CommentTok{{-}{-} décrire une table (ex. : \textbackslash{}d film)}
\NormalTok{\textbackslash{}du                }\CommentTok{{-}{-} lister les roles/users}
\NormalTok{\textbackslash{}timing            }\CommentTok{{-}{-} activer/désactiver le chronométrage des requêtes}
\NormalTok{\textbackslash{}i path.sql        }\CommentTok{{-}{-} exécuter les commandes depuis un fichier}
\NormalTok{\textbackslash{}?                 }\CommentTok{{-}{-} aide pour les meta{-}commands}
\end{Highlighting}
\end{Shaded}

\begin{center}\rule{0.5\linewidth}{0.5pt}\end{center}

\section{Vos Premières Requêtes
(Pagila)}\label{vos-premiuxe8res-requuxeates-pagila}

Connectez-vous à \texttt{pagila} :

\begin{Shaded}
\begin{Highlighting}[]
\NormalTok{\textbackslash{}c pagila}
\end{Highlighting}
\end{Shaded}

\subsection{Requêtes d'échauffement
(simples)}\label{requuxeates-duxe9chauffement-simples}

\begin{Shaded}
\begin{Highlighting}[]
\CommentTok{{-}{-} Combien de films ?}
\KeywordTok{SELECT} \FunctionTok{COUNT}\NormalTok{(}\OperatorTok{*}\NormalTok{) }\KeywordTok{FROM}\NormalTok{ film;}

\CommentTok{{-}{-} Voir quelques films}
\KeywordTok{SELECT} \OperatorTok{*} \KeywordTok{FROM}\NormalTok{ film }\KeywordTok{LIMIT} \DecValTok{5}\NormalTok{;}

\CommentTok{{-}{-} Films classés G}
\KeywordTok{SELECT}\NormalTok{ title, rating }\KeywordTok{FROM}\NormalTok{ film }\KeywordTok{WHERE}\NormalTok{ rating }\OperatorTok{=} \StringTok{\textquotesingle{}G\textquotesingle{}} \KeywordTok{LIMIT} \DecValTok{10}\NormalTok{;}
\end{Highlighting}
\end{Shaded}

\subsection{Requêtes
intermédiaires}\label{requuxeates-intermuxe9diaires}

\begin{Shaded}
\begin{Highlighting}[]
\CommentTok{{-}{-} Top 10 des films PG les plus chers à louer}
\KeywordTok{SELECT}\NormalTok{ title, rental\_rate}
\KeywordTok{FROM}\NormalTok{ film}
\KeywordTok{WHERE}\NormalTok{ rating }\OperatorTok{=} \StringTok{\textquotesingle{}PG\textquotesingle{}}
\KeywordTok{ORDER} \KeywordTok{BY}\NormalTok{ rental\_rate }\KeywordTok{DESC}\NormalTok{, title}
\KeywordTok{LIMIT} \DecValTok{10}\NormalTok{;}
\end{Highlighting}
\end{Shaded}

\subsection{Requêtes avancées pour
explorer}\label{requuxeates-avancuxe9es-pour-explorer}

\begin{Shaded}
\begin{Highlighting}[]
\CommentTok{{-}{-} Locations par client (top 10)}
\KeywordTok{SELECT}\NormalTok{ c.customer\_id,}
\NormalTok{       c.first\_name }\OperatorTok{||} \StringTok{\textquotesingle{} \textquotesingle{}} \OperatorTok{||}\NormalTok{ c.last\_name }\KeywordTok{AS}\NormalTok{ customer,}
       \FunctionTok{COUNT}\NormalTok{(}\OperatorTok{*}\NormalTok{) }\KeywordTok{AS}\NormalTok{ rentals}
\KeywordTok{FROM}\NormalTok{ rental r}
\KeywordTok{JOIN}\NormalTok{ customer c }\KeywordTok{USING}\NormalTok{ (customer\_id)}
\KeywordTok{GROUP} \KeywordTok{BY}\NormalTok{ c.customer\_id, customer}
\KeywordTok{ORDER} \KeywordTok{BY}\NormalTok{ rentals }\KeywordTok{DESC}
\KeywordTok{LIMIT} \DecValTok{10}\NormalTok{;}

\CommentTok{{-}{-} Revenus par catégorie}
\KeywordTok{SELECT}\NormalTok{ ca.name }\KeywordTok{AS} \KeywordTok{category}\NormalTok{,}
       \FunctionTok{SUM}\NormalTok{(p.amount) }\KeywordTok{AS}\NormalTok{ total\_revenue}
\KeywordTok{FROM}\NormalTok{ payment p}
\KeywordTok{JOIN}\NormalTok{ rental r }\KeywordTok{USING}\NormalTok{ (rental\_id)}
\KeywordTok{JOIN}\NormalTok{ inventory i }\KeywordTok{USING}\NormalTok{ (inventory\_id)}
\KeywordTok{JOIN}\NormalTok{ film f }\KeywordTok{USING}\NormalTok{ (film\_id)}
\KeywordTok{JOIN}\NormalTok{ film\_category fc }\KeywordTok{USING}\NormalTok{ (film\_id)}
\KeywordTok{JOIN} \KeywordTok{category}\NormalTok{ ca }\KeywordTok{USING}\NormalTok{ (category\_id)}
\KeywordTok{GROUP} \KeywordTok{BY}\NormalTok{ ca.name}
\KeywordTok{ORDER} \KeywordTok{BY}\NormalTok{ total\_revenue }\KeywordTok{DESC}\NormalTok{;}
\end{Highlighting}
\end{Shaded}

\begin{center}\rule{0.5\linewidth}{0.5pt}\end{center}

\section{Dépannage (Référence
Rapide)}\label{duxe9pannage-ruxe9fuxe9rence-rapide}

\begin{verbatim}
"psql n'est pas reconnu ..."
  - Utilisez Démarrer → SQL Shell (psql) (pas besoin de PATH), ou ajoutez
    C:\Program Files\PostgreSQL\<version>\bin au PATH ; rouvrez le terminal.

"password authentication failed for user 'postgres'"
  - Retapez soigneusement (sensible à la casse). Si vous pouvez vous connecter 
    via pgAdmin, changez le mot de passe là-bas. Complètement bloqué sur une 
    nouvelle installation ? La réinstallation peut être la plus rapide.

"could not connect to server" / timeouts
  - Assurez-vous que le service PostgreSQL fonctionne (Services). Vérifiez 
    host localhost et port 5432.

"Port already in use"
  - Une autre application utilise 5432. Arrêtez-la, ou installez PostgreSQL 
    sur un port différent (ex. : 5433) et utilisez ce port dans les clients.
\end{verbatim}

\begin{center}\rule{0.5\linewidth}{0.5pt}\end{center}

\section{Bonnes Pratiques de
Sécurité}\label{bonnes-pratiques-de-suxe9curituxe9}

\subsection{Créer un user non-superuser pour la
pratique}\label{cruxe9er-un-user-non-superuser-pour-la-pratique}

Au lieu d'utiliser toujours \texttt{postgres}, créez un user normal :

\begin{Shaded}
\begin{Highlighting}[]
\KeywordTok{CREATE} \KeywordTok{ROLE}\NormalTok{ student }\KeywordTok{WITH}\NormalTok{ LOGIN }\KeywordTok{PASSWORD} \StringTok{\textquotesingle{}mot\_de\_passe\_securise\textquotesingle{}}\NormalTok{;}
\KeywordTok{GRANT} \KeywordTok{ALL} \KeywordTok{PRIVILEGES} \KeywordTok{ON} \KeywordTok{DATABASE}\NormalTok{ pagila }\KeywordTok{TO}\NormalTok{ student;}

\CommentTok{{-}{-} Connectez{-}vous ensuite avec ce user}
\NormalTok{\textbackslash{}c pagila student}
\end{Highlighting}
\end{Shaded}

\textbf{Pourquoi ?} Le principe du moindre privilège : utilisez
\texttt{postgres} uniquement pour l'administration, pas pour les
requêtes quotidiennes.

\begin{center}\rule{0.5\linewidth}{0.5pt}\end{center}

\section{Erreurs Courantes des
Étudiants}\label{erreurs-courantes-des-uxe9tudiants}

\begin{itemize}
\tightlist
\item
  ❌ \textbf{Oublier les points-virgules} dans les instructions SQL
\item
  ❌ \textbf{Utiliser des backslashes Windows}
  (\texttt{\textbackslash{}}) dans les chemins de fichiers dans
  \texttt{psql} (utilisez \texttt{/})
\item
  ❌ \textbf{Rester connecté à la mauvaise database} (utilisez
  \texttt{\textbackslash{}c} pour changer)
\item
  ❌ \textbf{Exécuter des commandes DDL sans réfléchir} (DROP, DELETE
  sans WHERE)
\item
  ❌ \textbf{Mélanger SQL et meta-commands} (les meta-commands n'ont pas
  besoin de \texttt{;})
\end{itemize}

\begin{center}\rule{0.5\linewidth}{0.5pt}\end{center}

\section{Et Maintenant ? Prochaines
Étapes}\label{et-maintenant-prochaines-uxe9tapes}

\subsection{Concepts importants à apprendre
ensuite}\label{concepts-importants-uxe0-apprendre-ensuite}

\begin{itemize}
\tightlist
\item
  \textbf{Indexes} --- accélérer les requêtes
\item
  \textbf{Transactions} --- garantir la cohérence des données
\item
  \textbf{Constraints} --- clés primaires, clés étrangères, checks
\item
  \textbf{Views} --- requêtes enregistrées
\item
  \textbf{Functions et Procedures} --- logique réutilisable
\end{itemize}

\subsection{Ressources recommandées}\label{ressources-recommanduxe9es}

\begin{itemize}
\tightlist
\item
  \textbf{Documentation officielle PostgreSQL :}
  \url{https://www.postgresql.org/docs/}
\item
  \textbf{PostgreSQL Exercises (pratique interactive) :}
  \url{https://pgexercises.com/}
\item
  \textbf{Tutoriel PostgreSQL :}
  \url{https://www.postgresqltutorial.com/}
\item
  \textbf{Pagila GitHub (plus d'exemples) :}
  \url{https://github.com/devrimgunduz/pagila}
\end{itemize}

\subsection{Exercices pratiques avec
Pagila}\label{exercices-pratiques-avec-pagila}

\begin{enumerate}
\def\labelenumi{\arabic{enumi}.}
\tightlist
\item
  Trouvez tous les films avec l'acteur ``PENELOPE GUINESS''
\item
  Calculez la durée moyenne de location par catégorie de film
\item
  Identifiez les clients qui n'ont jamais loué de film
\item
  Créez une view montrant les films en stock dans chaque magasin
\item
  Écrivez une requête pour trouver les films les plus rentables
\end{enumerate}

\subsection{Avant de demander de l'aide ---
Checklist}\label{avant-de-demander-de-laide-checklist}

\begin{itemize}
\tightlist
\item
  ✅ Le service PostgreSQL fonctionne-t-il ?
\item
  ✅ Avez-vous essayé de vous connecter avec \texttt{psql} ?
\item
  ✅ Êtes-vous connecté à la bonne database ?
\item
  ✅ Avez-vous vérifié l'orthographe (tables, colonnes, mots de passe) ?
\item
  ✅ Avez-vous consulté les messages d'erreur complets ?
\item
  ✅ Avez-vous cherché l'erreur dans la documentation ou en ligne ?
\end{itemize}

\begin{center}\rule{0.5\linewidth}{0.5pt}\end{center}

\section{Sauvegarde et Restauration
(Bases)}\label{sauvegarde-et-restauration-bases}

\subsection{Sauvegarder une database}\label{sauvegarder-une-database}

\textbf{CMD :}

\begin{Shaded}
\begin{Highlighting}[]
\BuiltInTok{set} \VariableTok{PGBIN}\NormalTok{=C:\textbackslash{}Program Files\textbackslash{}PostgreSQL\textbackslash{}17\textbackslash{}bin}
\StringTok{"}\PreprocessorTok{\%}\VariableTok{PGBIN}\PreprocessorTok{\%}\StringTok{\textbackslash{}pg\_dump.exe"}\NormalTok{ {-}U postgres }\AttributeTok{{-}d}\NormalTok{ pagila }\AttributeTok{{-}f}\NormalTok{ C:\textbackslash{}backup\_pagila.sql}
\end{Highlighting}
\end{Shaded}

\textbf{PowerShell :}

\begin{Shaded}
\begin{Highlighting}[]
\VariableTok{$PGBIN} \OperatorTok{=} \StringTok{"C:\textbackslash{}Program Files\textbackslash{}PostgreSQL\textbackslash{}17\textbackslash{}bin"}
\OperatorTok{\&} \StringTok{"}\VariableTok{$PGBIN}\StringTok{\textbackslash{}pg\_dump.exe"} \OperatorTok{{-}}\NormalTok{U postgres }\OperatorTok{{-}}\NormalTok{d pagila }\OperatorTok{{-}}\NormalTok{f C}\OperatorTok{:}\NormalTok{\textbackslash{}backup\_pagila}\OperatorTok{.}\FunctionTok{sql}
\end{Highlighting}
\end{Shaded}

\subsection{Restaurer une database}\label{restaurer-une-database}

\textbf{CMD :}

\begin{Shaded}
\begin{Highlighting}[]
\StringTok{"}\PreprocessorTok{\%}\VariableTok{PGBIN}\PreprocessorTok{\%}\StringTok{\textbackslash{}psql.exe"}\NormalTok{ {-}U postgres }\AttributeTok{{-}d}\NormalTok{ pagila }\AttributeTok{{-}f}\NormalTok{ C:\textbackslash{}backup\_pagila.sql}
\end{Highlighting}
\end{Shaded}

\textbf{PowerShell :}

\begin{Shaded}
\begin{Highlighting}[]
\OperatorTok{\&} \StringTok{"}\VariableTok{$PGBIN}\StringTok{\textbackslash{}psql.exe"} \OperatorTok{{-}}\NormalTok{U postgres }\OperatorTok{{-}}\NormalTok{d pagila }\OperatorTok{{-}}\NormalTok{f C}\OperatorTok{:}\NormalTok{\textbackslash{}backup\_pagila}\OperatorTok{.}\FunctionTok{sql}
\end{Highlighting}
\end{Shaded}

\textbf{Conseil :} Sauvegardez régulièrement vos databases de pratique
avant d'expérimenter des commandes destructrices !

\begin{center}\rule{0.5\linewidth}{0.5pt}\end{center}

\section{Vous Êtes Prêt !}\label{vous-uxeates-pruxeat}

Pratiquez un peu chaque jour --- avec \texttt{psql} ou votre GUI préféré
--- et vous vous sentirez rapidement à l'aise. Explorez les schemas et
les clés, essayez \texttt{EXPLAIN\ ANALYZE}, et gardez des notes des
commandes que vous exécutez. Bon apprentissage et bonnes requêtes ! 🎉




\end{document}
